% test_full.tex - 完整模板测试
% 测试所有主要功能：封面、摘要、目录、正文、参考文献、致谢

\documentclass[degree=bachelor]{sduthesis}

% 配置文件
\SDUSetup{
  title = {基于深度学习的图像识别方法研究},
  author = {张三},
  student-id = {2021001234},
  school = {计算机科学与技术学院},
  major = {计算机科学与技术},
  supervisor = {李教授},
  supervisor-title = {教授},
  date = {2025}年6月
}

\begin{document}

% 封面
\MakeCover

% 中文摘要
\begin{cnabstract}
本文研究了基于深度学习的图像识别方法，主要贡献包括：
\begin{itemize}
  \item 提出了一种改进的卷积神经网络结构
  \item 设计了新的特征提取模块
  \item 在多个数据集上验证了方法的有效性
\end{itemize}
关键词：深度学习，图像识别，卷积神经网络，特征提取
\end{cnabstract}

% 英文摘要
\begin{enabstract}
This thesis investigates image recognition methods based on deep learning. The main contributions include:
\begin{itemize}
  \item Proposing an improved convolutional neural network architecture
  \item Designing a novel feature extraction module
  \item Validating the effectiveness on multiple datasets
\end{itemize}
Keywords: Deep Learning, Image Recognition, CNN, Feature Extraction
\end{enabstract}

% 目录
\tableofcontents

% 正文
\chapter{引言}
\label{chapter:intro}

\section{研究背景}
图像识别是计算机视觉领域的重要研究方向，近年来深度学习技术的
发展极大地推动了图像识别性能的提升~\cite{lecun2015deep}。

\section{研究意义}
深度学习方法在图像识别任务上取得了显著成果，但仍存在一些
挑战需要解决。

\chapter{相关工作}
\label{chapter:related}

\section{传统方法}
传统的图像识别方法主要依赖手工设计的特征，如 SIFT~\cite{lowe2004distinctive}、
HOG~\cite{dalal2005histograms} 等。

\section{深度学习方法}
卷积神经网络（CNN）~\cite{krizhevsky2012imagenet}是深度学习在图像识别
领域最成功的应用之一。

\chapter{方法}
\label{chapter:method}

\section{网络结构}
本文提出的网络结构如图~\ref{fig:network}所示。

\begin{figure}[htbp]
  \centering
  \includegraphics[width=0.8\textwidth]{example-image}
  \caption{网络结构示意图}
  \label{fig:network}
\end{figure}

\section{损失函数}
模型使用交叉熵损失函数：
\begin{equation}
  \mathcal{L} = -\sum_{i} y_i \log(\hat{y}_i)
  \label{eq:loss}
\end{equation}

\subsection{优化方法}
使用 Adam 优化器进行训练，学习率设为 0.001。

\chapter{实验}
\label{chapter:experiment}

\section{数据集}
实验在 ImageNet~\cite{deng2009imagenet} 数据集上进行。

\section{结果分析}
实验结果如表~\ref{tab:results}所示。

\begin{table}[htbp]
  \centering
  \caption{实验结果对比}
  \begin{tabular}{ccc}
    \toprule
    方法 & 准确率 & 参数量 \\
    \midrule
    ResNet-50 & 76.1\% & 25.6M \\
    本文方法 & 78.3\% & 23.1M \\
    \bottomrule
  \end{tabular}
  \label{tab:results}
\end{table}

\chapter{结论}
\label{chapter:conclusion}

本文提出了一种基于深度学习的图像识别方法，通过改进网络结构和
特征提取模块，取得了较好的实验效果。未来工作将探索更复杂的
网络结构和更大规模的数据集。

% 参考文献
\nocite{*}
\bibliography{test_bib}

% 致谢
\begin{acknowledgement}
感谢导师李教授的悉心指导，感谢实验室同学的帮助与支持。
\end{acknowledgement}

\end{document}
